%==============================================================================
% Independent Verification of a 29 998-Digit Pocklington–Lehmer Certificate
% H. Bakkaoui — May 2026
%==============================================================================
\documentclass[11pt,a4paper]{article}

\usepackage[utf8]{inputenc}
\usepackage[T1]{fontenc}
\usepackage[english]{babel}
\usepackage{amsmath,amssymb,amsthm,mathtools}
\usepackage{geometry}
\geometry{margin=2.5cm}
\usepackage{booktabs}
\usepackage{array}
\usepackage{enumitem}
\usepackage[dvipsnames]{xcolor}
\usepackage{tcolorbox}
\tcbuselibrary{skins,breakable}
\usepackage{hyperref}
\hypersetup{colorlinks=true,linkcolor=MidnightBlue,
            citecolor=MidnightBlue,urlcolor=MidnightBlue}

%--- Theorem environments -----------------------------------------------------
\newtheorem{theorem}{Theorem}
\newtheorem{lemma}[theorem]{Lemma}
\theoremstyle{definition}
\newtheorem{definition}{Definition}
\theoremstyle{remark}
\newtheorem{remark}{Remark}

%--- Coloured boxes -----------------------------------------------------------
\newtcolorbox{rigorousbox}[1][]{
  colback=Green!8, colframe=Green!60!black, arc=1mm,
  title=\textbf{Rigorous}, fonttitle=\bfseries, breakable, #1}

\newtcolorbox{certbox}[1][]{
  colback=RoyalBlue!5, colframe=RoyalBlue!70!black, arc=1mm,
  title=\textbf{Certificate}, fonttitle=\bfseries, breakable, #1}

\newtcolorbox{resultbox}[1][]{
  colback=Goldenrod!8, colframe=Goldenrod!70!black, arc=1mm,
  title=\textbf{Verification Record}, fonttitle=\bfseries, breakable, #1}

%--- Shortcuts ----------------------------------------------------------------
\newcommand{\N}{\mathbb{N}}

%==============================================================================
\title{Independent Verification of a 29\,998-Digit\\
       Pocklington--Lehmer Primality Certificate\\[6pt]
       {\large Parametric family $p = 3m(m+1)+1$,\;
               $m = 2^a\cdot 3^b - 1$}}

\author{Hassane \textsc{Bakkaoui}\\
\small Independent Researcher\\
\small \texttt{bakkahassa@hotmail.com}}

\date{May 2026}
%==============================================================================

\begin{document}
\maketitle
\thispagestyle{empty}

\begin{abstract}
We report the complete, independently verified Pocklington--Lehmer primality
certificate for
\[
  p \;=\; 3m(m+1)+1, \qquad m = 2^{19435}\cdot 3^{19173}-1,
\]
a prime of $\mathbf{29\,998}$ \textbf{decimal digits}.
The certificate was originally produced by the \textsc{PrimeQuest v4}
algorithm in $3\,\text{h}\,06\,\text{min}$ on a single consumer laptop.
Here we verify it independently using exact arbitrary-precision arithmetic
(Python~3.12 / GMP), with \emph{no probabilistic step whatsoever}:
all four Pocklington--Lehmer conditions are confirmed, establishing the
unconditional primality of~$p$.
\end{abstract}

\tableofcontents
\newpage

%==============================================================================
\section{The Number and Its Parameters}
%==============================================================================

\begin{definition}[29\,998-digit PrimeQuest instance]
\label{def:instance}
Set
\[
  a = 19\,435, \qquad b = 19\,173,
\]
and define
\begin{align*}
  m &\;=\; 2^{19435}\cdot 3^{19173} - 1
       \qquad (14\,999 \text{ decimal digits}),\\
  p &\;=\; 3m(m+1)+1
       \qquad\;\; (29\,998 \text{ decimal digits}).
\end{align*}
\end{definition}

\subsection*{Identification --- first and last digits of $p$}

The decimal expansion of $p$ begins and ends as follows:
\begin{equation}
\label{eq:pdigits}
  p \;=\;
  \underbrace{1602087359509060348500037851549319056715}_{\text{first 40 digits}}
  \;\cdots\;
  \underbrace{16930852697403293697}_{\text{last 20 digits}}.
\end{equation}
These values were reconstructed from scratch during the independent
verification (Step~1 of Section~\ref{sec:verif}) and match the stored
certificate exactly.

\subsection*{Digit-count formula}

The number of decimal digits of $p(a,b) = 3m(m+1)+1$ is
\[
  D(a,b)
  \;=\; \bigl\lfloor 2a\log_{10}2 + (2b+1)\log_{10}3 \bigr\rfloor + 1.
\]
For $(a,b)=(19435,19173)$:
\[
  D \;=\; \lfloor 2\times 19435\times 0.30103
          + 39347\times 0.47712 \rfloor + 1
  \;=\; 29\,998. \checkmark
\]

%==============================================================================
\section{Factorisation of $p-1$ and the Pocklington Bound}
%==============================================================================

\begin{rigorousbox}
\begin{lemma}[Explicit factorisation of $p-1$]
\label{lem:factorisation}
With $m$ and $p$ as in Definition~\ref{def:instance},
\[
  p - 1 \;=\;
  \underbrace{2^{19435}\cdot 3^{19174}}_{F} \;\cdot\; m,
\]
where $F := 2^{19435}\cdot 3^{19174}$ has exactly $15\,000$ decimal digits
and satisfies $F > \sqrt{p}$.
\end{lemma}

\begin{proof}
Since $m + 1 = 2^{19435}\cdot 3^{19173}$,
\[
  p - 1 \;=\; 3m(m+1)
           \;=\; 3m\cdot 2^{19435}\cdot 3^{19173}
           \;=\; 2^{19435}\cdot 3^{19174}\cdot m.
\]
Because $p = 3m^2+3m+1 < 3(m+1)^2$, we get
$\sqrt{p} < \sqrt{3}\,(m+1) < 3(m+1) = 3\cdot 2^{19435}\cdot 3^{19173}
= 2^{19435}\cdot 3^{19174} = F$.
\end{proof}
\end{rigorousbox}

\subsection*{Numerical confirmation of $F > \sqrt{p}$}

Comparing base-2 logarithms:
\[
  2\log_2 F
  \;=\; 2\bigl(19435 + 19174\log_2 3\bigr)
  \;=\; 99\,650.1420,
\]
\[
  \log_2 p \;=\; 99\,648.5570.
\]

\begin{center}
\renewcommand{\arraystretch}{1.3}
\begin{tabular}{@{}lc@{}}
\toprule
Quantity & Value \\
\midrule
$2\log_2 F$ & $99\,650.1420$ \\
$\log_2 p$  & $99\,648.5570$ \\
Margin $= 2\log_2 F - \log_2 p$ & $\mathbf{+1.585}$ bits \\
\midrule
$F > \sqrt{p}$ & \textbf{True} \checkmark \\
\bottomrule
\end{tabular}
\end{center}

\noindent
The margin of $1.585$ bits means $F \approx 3\sqrt{p}$, well above the
Pocklington threshold.

%==============================================================================
\section{The Pocklington--Lehmer Certificate}
%==============================================================================

\subsection*{Criterion}

The Pocklington--Lehmer theorem \cite{Pocklington1914,Lehmer1927} states:

\begin{theorem}[Pocklington--Lehmer]
\label{thm:pocklington}
Let $N > 1$ be an integer. Suppose there exists $F \mid N-1$ with $F > \sqrt{N}$,
and that for each prime $q \mid F$ there exists an integer $w_q$ such that
\[
  w_q^{N-1} \equiv 1 \pmod{N}
  \quad\text{and}\quad
  \gcd\!\bigl(w_q^{(N-1)/q} - 1,\; N\bigr) = 1.
\]
Then $N$ is prime.
\end{theorem}

\noindent
For $p$ as in Definition~\ref{def:instance}, Lemma~\ref{lem:factorisation}
provides $F = 2^{19435}\cdot 3^{19174}$ with $F > \sqrt{p}$. The prime
divisors of $F$ are $q = 2$ and $q = 3$ only, so \emph{exactly two witnesses
are needed}.

\subsection*{The certificate}

\begin{certbox}[title={\textbf{Pocklington--Lehmer Certificate} ---
                        29\,998-digit prime}]
\[
  p \;=\; 3m(m+1)+1, \qquad m = 2^{19435}\cdot 3^{19173}-1.
\]

\medskip
\noindent\textbf{Factored part of $p-1$:}\quad
$F = 2^{19435}\cdot 3^{19174}, \quad F > \sqrt{p}$. \checkmark

\bigskip
\renewcommand{\arraystretch}{1.6}
\begin{tabular}{@{}ccccc@{}}
\toprule
Prime $q$ & Witness $w_q$
  & $w_q^{p-1}\!\equiv 1\pmod{p}$
  & $\gcd(w_q^{(p-1)/q}\!-1,p)=1$ & Status \\
\midrule
$2$ & $5$ & \checkmark & \checkmark & \textbf{OK} \\
$3$ & $7$ & \checkmark & \checkmark & \textbf{OK} \\
\bottomrule
\end{tabular}

\bigskip
\noindent\textbf{Conclusion:} $p$ is \emph{unconditionally prime}
(deterministic proof). \checkmark
\end{certbox}

%==============================================================================
\section{Independent Numerical Verification}
\label{sec:verif}
%==============================================================================

All four Pocklington--Lehmer conditions were re-checked from scratch
using \texttt{Python~3.12} with its built-in GMP-backed integer arithmetic.
No probabilistic step was used at any point.

\begin{resultbox}[title={\textbf{Verification Record} --- independent computation,
                          May 2026}]

\noindent\textbf{Environment:}
\begin{itemize}[nosep,leftmargin=1.5em]
  \item Python 3.12 (CPython), GMP backend.
  \item \texttt{sys.set\_int\_max\_str\_digits(40\,000)} to allow conversion
        of 30\,000-digit integers.
  \item All exponentiations: built-in \texttt{pow(base, exp, mod)} --- exact
        modular arithmetic, no approximation.
  \item Hardware: single CPU core, no multi-threading.
\end{itemize}

\bigskip
\noindent\textbf{Step 1 --- Reconstruct $p$ from parameters.}\\
Computed $m = 2^{19435}\cdot 3^{19173}-1$ and $p = 3m(m+1)+1$ from $(a,b)$.

\medskip
\begin{center}
\renewcommand{\arraystretch}{1.3}
\begin{tabular}{@{}lll@{}}
\toprule
 & Computed & Stored certificate \\
\midrule
Digits of $m$ & $14\,999$ & $14\,999$ \\
Digits of $p$ & $29\,998$ & $29\,998$ \\
First 40 digits of $p$ & \texttt{16020873595090603485\ldots056715}
                        & \texttt{16020873595090603485\ldots056715} \\
Last 20 digits of $p$  & \texttt{16930852697403293697}
                        & \texttt{16930852697403293697} \\
\midrule
Match & \multicolumn{2}{c}{\textbf{Exact} \checkmark} \\
\bottomrule
\end{tabular}
\end{center}

\bigskip
\noindent\textbf{Step 2 --- Pocklington bound.}
\[
  2\log_2 F = 99\,650.1420 \;>\; \log_2 p = 99\,648.5570. \checkmark
\]

\bigskip
\noindent\textbf{Step 3 --- Witness $w_2 = 5$ for $q = 2$.}
\begin{align*}
  5^{p-1} &\equiv 1 \pmod{p}. \checkmark \\
  \gcd\!\bigl(5^{(p-1)/2}-1,\;p\bigr) &= 1. \checkmark
\end{align*}

\noindent\textbf{Step 4 --- Witness $w_3 = 7$ for $q = 3$.}
\begin{align*}
  7^{p-1} &\equiv 1 \pmod{p}. \checkmark \\
  \gcd\!\bigl(7^{(p-1)/3}-1,\;p\bigr) &= 1. \checkmark
\end{align*}

\bigskip
\noindent\textbf{Timing.}

\medskip
\begin{center}
\renewcommand{\arraystretch}{1.3}
\begin{tabular}{@{}ll@{}}
\toprule
Operation & Wall-clock time \\
\midrule
Reconstruct $m$, $p$ & $< 1$ s \\
$5^{p-1} \bmod p$            & $1\,735.9$ s \\
$\gcd(5^{(p-1)/2}-1,\,p)$    & $\phantom{1\,73}0.6$ s \\
$7^{p-1} \bmod p$            & $1\,736.0$ s \\
$\gcd(7^{(p-1)/3}-1,\,p)$    & $\phantom{1\,73}0.6$ s \\
\midrule
\textbf{Total} & $\mathbf{6\,945}$ \textbf{s}
                 $\approx \mathbf{1\,\text{h}\;56\,\text{min}}$ \\
\bottomrule
\end{tabular}
\end{center}

\bigskip
\noindent\textbf{Conclusion:}
All four Pocklington--Lehmer conditions are satisfied.
The integer $p$ is \textbf{unconditionally prime}.
\end{resultbox}

%==============================================================================
\section{Comparison: Original Search vs.\ Independent Verification}
%==============================================================================

\begin{center}
\renewcommand{\arraystretch}{1.3}
\begin{tabular}{@{}lll@{}}
\toprule
 & \textbf{PrimeQuest v4 (search)} & \textbf{This work (verification)} \\
\midrule
Goal      & Find a prime candidate         & Check the certificate \\
Algorithm & Sieve + MR pre-filter + Pocklington
           & Pure Pocklington--Lehmer \\
Probabilistic steps & 286 Miller--Rabin tests       & \textbf{None} \\
Hardware  & ARM64, 9 cores (Snapdragon X Plus) & Single CPU core \\
Wall-clock & $11\,146$ s ($3$ h $06$ min)   & $6\,945$ s ($1$ h $56$ min) \\
Speed-up factor & $1\times$ (reference)        & $\approx 1.6\times$ faster \\
Result    & Certificate produced           & Certificate \textbf{confirmed} \\
\bottomrule
\end{tabular}
\end{center}

\medskip
\noindent
The verification is faster because it skips the sieve, the Theorem-2
algebraic filter (which eliminates $33\%$ of candidates), and the 286
Miller--Rabin pre-screening rounds. It runs exactly four modular
exponentiations on the already-known 29\,998-digit modulus.

%==============================================================================
\section{Summary}
%==============================================================================

\begin{rigorousbox}[title={\textbf{Final statement}}]
The integer
\[
  p \;=\; 3m(m+1)+1, \qquad m = 2^{19435}\cdot 3^{19173}-1,
\]
with $29\,998$ decimal digits, beginning
\[
  p = 1602087359509060348500037851549319056715\cdots 16930852697403293697,
\]
is \textbf{prime}.

\medskip
\noindent
The unconditional proof is provided by the Pocklington--Lehmer criterion
(Theorem~\ref{thm:pocklington}):
\begin{itemize}[nosep]
  \item $F = 2^{19435}\cdot 3^{19174} > \sqrt{p}$ \hfill\checkmark
  \item $w_2 = 5$ is a Pocklington witness for $q=2$ \hfill\checkmark
  \item $w_3 = 7$ is a Pocklington witness for $q=3$ \hfill\checkmark
\end{itemize}

\medskip
\noindent
All conditions were verified by an independent exact computation
(Python~3.12 / GMP, $6\,945$ s, single core, no probabilistic step).
\end{rigorousbox}

%==============================================================================
\begin{thebibliography}{9}

\bibitem{Pocklington1914}
H.\,C.~Pocklington,
\textit{The determination of the prime or composite nature of large numbers
by Fermat's theorem},
Proc.\ Cambridge Phil.\ Soc.\ \textbf{18} (1914--16), 29--30.

\bibitem{Lehmer1927}
D.\,H.~Lehmer,
\textit{Tests for primality by the converse of Fermat's theorem},
Bull.\ Amer.\ Math.\ Soc.\ \textbf{33} (1927), 327--340.

\bibitem{BakkPaper4}
H.~Bakkaoui,
\textit{Unconditional Primality Certificates for Numbers of the Form
$p = 3m(m+1)+1$: Arithmetic Filters, Multi-Core Search, and Records to
$29\,998$ Decimal Digits},
preprint, 2026.

\end{thebibliography}

\end{document}
